\documentclass[a4paper,10pt]{report}\input{../0.Préambule/Préambule_cours_prof}\begin{document}

\definecolor{titlecolor}{rgb}{0.98, 0.4, 0.0}
\newpage
\setcounter{chapter}{9}
\chapter{Nombres relatifs 2}

\section{Somme de deux nombres relatifs}

\vspace{-3mm}
\subsection{Cas de deux nombres relatifs de même signe}

\begin{prop}
\begin{enumerate}
\item on additionne les distances à zéro
\item on donne au résultat le signe commun aux deux nombres.
\end{enumerate}
\end{prop}

\begin{ex}
\noindent \renewcommand{\arraystretch}{1}
\begin{tabularx}{\linewidth}{*{4}{>{\arraybackslash}X}}
$(- 4) + (- 2) = ( - 6)$ &
$(+3) + (+7) = (+10)$&
$(+7,8) + (+8,7)= (+ 16,5)$ &
$(-11,2) + (- 6,4) = (-17,6)$\\
\end{tabularx}
\end{ex}

\vspace{-3mm}
\subsection{Cas de deux nombres relatifs de signes contraires}

\begin{prop}
\begin{enumerate}
\item on soustrait les distances à zéro
\item on donne au résultat le signe du nombre qui a la plus grande distance à zéro.
\end{enumerate}
\end{prop}

\begin{ex}
\noindent \renewcommand{\arraystretch}{1}
\begin{tabularx}{\linewidth}{*{4}{>{\arraybackslash}X}}
$(+ 3) + (-7) = (- 4)$ &
$(-5) + (+3) = (- 2)$&
$(+7,8) + (- 4,5) = (+3,3)$ &
$4 + (- 8) = (- 4)$ \\
\end{tabularx}
\end{ex}

\begin{rmq}
La somme de deux nombres relatifs opposés est égale à zéro
\end{rmq}

\begin{ex}
\noindent \renewcommand{\arraystretch}{1}
\begin{tabularx}{\linewidth}{*{2}{>{\arraybackslash}X}}
$(+3,4) + (- 3,4) = 0$&
$(- 7) + (+ 7) = 0$\\
\end{tabularx}
\end{ex}

\vspace{-3mm}
\subsection{Somme de plusieurs termes}

\begin{prop}
Pour calculer la somme de plusieurs termes, on peut :
\begin{enumerate}
\item soit calculer dans l'ordre de lecture, de gauche à droite
\item soit modifier l'ordre et les regrouper par signe.
\end{enumerate}
\end{prop}

\begin{ex}
\vspace{-6mm}
\begin{minipage}{7cm}
\begin{eqnarray*}
A &=& (- 3,5) + (+6) + (- 8)\\ 
A &=& (- 3,5) + (- 8) + (+ 6)\\
A &=& (- 11,5) + (+6)\\
A &=& (- 5,5)
\end{eqnarray*}
\end{minipage}
\begin{minipage}{7cm}
\begin{eqnarray*}
B &=& (- 2,3) + (+ 4,3) + (- 9) + (+ 2,3) + (+3,8) + (- 5,7) \\
&&\text{Il y a deux opposés : } (- 2,3) \text{ et } (+ 2,3)\\
&=& (+ 4,3) + (+ 3,8) + (- 9) + (- 5,7)\\
&=& (+8,1) + (- 14,7)\\
&=& (- 6,6)
\end{eqnarray*}
\end{minipage}
\end{ex}

\vspace{-6mm}
\section{Soustraction}

\begin{prop}
Pour soustraire un nombre relatif, on ajoute son opposé.
\end{prop}

\begin{ex}
\noindent \renewcommand{\arraystretch}{1.5}
\begin{tabularx}{\linewidth}{*{2}{>{\arraybackslash}X}}
$(-5) - (+ 20) = (-5) + (-20) = (- 25)$ &
$(+7) - (-4) = (+7) + (+4) = (+11)$\\
$(+15) - (+18) = (+15) + (- 18) = (- 3)$ &
$(- 8) - (- 5) = (- 8) + (+ 5) = (- 3)$\\
\end{tabularx}
\end{ex}

\vspace{-3mm}
\section{Simplification d'écriture}

\subsection{Règles de simplification}

\begin{prop}
Pour écrire plus simplement une suite d'additions et de soustractions de nombres relatifs :
\begin{enumerate}
\item On transforme les soustractions en additions.
\item On supprime les signes + des additions et les parenthèses autour des nombres.
\item On supprime le signe + devant un nombre si il est en début de calcul.
\end{enumerate}
\end{prop}

\begin{ex}
\noindent \renewcommand{\arraystretch}{1.5}
\begin{tabularx}{\linewidth}{*{2}{>{\arraybackslash}X}}
$(-5) + (+3) = - 5 + 3 = -2$ &
$(+ 4) + (-2) = 4 - 2 = 2$ \\
$(- 7) - (+ 2) = (- 7) + (- 2) =-7 - 2 = -9$ &
$(- 5) - (- 4) = (- 5) + ( 4) = -5 + 4 = -1$\\
\end{tabularx}
\end{ex}

\begin{rmq}
\renewcommand{\arraystretch}{1.8}
\begin{tabularx}{\linewidth}{|*{3}{>{ \centering \arraybackslash }X|}}
\rowcolor{lightgray!50} Écriture simplifiée & Signification & Résultat \\\hline
 $6+9$ & $(+6) + (+9)$ & $ (+15) = 15$ \\\hline
 $7-13$ & $(+7) + (-13)$ & $ (-6) = -6$ \\\hline
 $-6,5+10$ & $(-6,5) + (+10)$ & $ (+3,5) =3,5$ \\\hline
 $-4,2-5,8$ & $(-4,2) + (-5,8)$ & $ (-10) = -10$ \\\hline
\end{tabularx}
\end{rmq}

\subsection{Somme algébrique}

\begin{dft}
Une somme algébrique est une suite d'additions et de soustractions.
\end{dft}

\begin{ex}
\noindent \begin{minipage}{8cm}
\begin{eqnarray*}
A &=& (-7) +(+5) - (-9 ) + (+2) - (+6) \\
&=& (-7) + (+5) + (+9) + (+2) + (-6) \\
&&\text{On transforme les soustractions en additions.}\\
&=& -7 + 5 + 9 + 2 - 6 \\
&&\text{On simplifie l'écriture.}\\
&=& - 7 - 6 + 5 + 9 + 2\\
&=& - 13 + 16\\
&=& 3
\end{eqnarray*}
\end{minipage}
\begin{minipage}{8cm}
\begin{eqnarray*}
B &=& - 4,5 + 6 + 4,5 - 3 + 4 - 2 \\
&& \text{l'écriture est déjà simplifiée ;}\\
&& \text{on regarde s'il y a des opposés :}\\
&=& 6 - 3 + 4 - 2\\
&& \text{on regroupe les nombres positifs entre eux} \\
&& \text{et les nombres négatifs entre eux}\\
&=& 6 + 4 - 3 - 2\\
&=& 10 - 5\\
&=& 5
\end{eqnarray*}
\end{minipage}
\end{ex}

\end{document}